\chapter{Introduction}

\par The transition from academic theory to professional engineering is often abrupt. While university courses provide a strong foundation in algorithms and syntax, they rarely simulate the volatility of a production environment. In a classroom, code is written to pass test cases; in the industry, code is written to survive user interaction, scalability demands, and market pivots. This industrial training at \textbf{Tanbel Inc.} was not just an application of skills—it was an immersion into the reality of modern software development. I moved from writing isolated functions to contributing to a living, breathing product ecosystem, where "urgent" means "now," and stability is the primary metric of success.

\par The objective of this training was to bridge the gap between student mindset and engineer mindset. It required navigating complex codebases, understanding the business logic behind B2C and B2B operations, and delivering features that directly impacted the company's bottom line.

\section{Problem Domain}

\par \textbf{Tanbel Inc.} operates at the intersection of logistics and waste management—an industry traditionally bogged down by manual processes and paper trails. The problem is coordination: connecting customers who need junk removed with haulers who have the capacity to remove it, instantly and reliably.

\par To solve this, Tanbel Inc. deploys two flagship products:
\begin{itemize}
    \item \textbf{WeCycle} (\href{https://getwecycle.com/}{getwecycle.com}): A B2C platform that allows homeowners and businesses to book junk removal services on-demand. Think of it as the "Uber" for junk.
    \item \textbf{ArrayQ} (\href{https://arrayq.com/}{arrayq.com}): A B2B CRM and dashboard solution designed for haulers to manage their dumpsters, inventory, and schedules efficiently.
\end{itemize}

\par My role as a \textbf{Full-Stack Developer} was to maintain and enhance the digital infrastructure that powers these solutions. Specifically, I worked on the WeCycle Dashboard Application (WCAPP) and the ArrayQ mobile/web platforms, ensuring that the connection between the customer’s request and the hauler’s execution was seamless.

\par As a trainee, I didn’t just watch; I had to get my hands dirty. I ran into a bunch of hurdles:
\begin{itemize}
    \item We had to rip out the old Ant Design components and replace them with our custom web-ui library. It was tedious, but necessary.
    \item I had to make the WCAPP look good on everything—phones, tablets, desktops. If it looked bad on an iPad, it was my problem.
    \item Debugging state management in a massive app is... painful. But I had to figure it out.
    \item I also had to learn the codebase without breaking it. It’s intimidating to push code to a repo that people actually depend on.
\end{itemize}

\par I couldn't have done it alone. My team—supervisors and colleagues—were a huge help. We used the modern stack: React.js, Next.js, and TypeScript. Honestly, it’s the standard now for a reason.

\section{Motivation of Training}

\par Why did I do this? A few reasons:

\begin{itemize}
    \item \textbf{Actually Using What I Learned:} I wanted to see if the software engineering theories I sat through in class actually hold up. (Spoiler: Some do, some don't.)
    \item \textbf{Getting Better, Faster:} I needed to get sharp with React.js and Next.js. TypeScript was another beast I wanted to tame.
    \item \textbf{Agile Chaos:} I wanted to see what "Agile" feels like when deadlines are real.
    \item \textbf{Meeting People:} It’s not just about code. I wanted to meet developers who have been doing this for years and pick their brains.
    \item \textbf{Future-Proofing:} I’m trying to build a career here. This was the groundwork.
\end{itemize}

\section{Training Goals}

\par My goals were pretty straightforward:

\begin{itemize}
    \item Get really, really good at React.js and Next.js.
    \item Stop solving textbook problems and start solving business problems.
\end{itemize}

\section{Scope of Training}

\section{Scope of Training}

\par My role was multifaceted, operating across the entire application stack. As a \textbf{Full-Stack Developer}, I was responsible for end-to-end feature delivery—from designing database schemas to deploying pixel-perfect UIs.

\par The scope included:
\begin{itemize}
    \item \textbf{Frontend Engineering:} Building responsive interfaces with React.js, Next.js, and Capacitor for mobile.
    \item \textbf{Backend Development:} Architecting scalable APIs using Node.js and Express.js, and managing data integrity with MongoDB and PostgreSQL.
    \item \textbf{DevOps \& Infrastructure:} Configuring Docker containers, managing Nginx servers, and optimizing CI/CD workflows with Nx and GitHub Actions.
    \item \textbf{System Integration:} Connecting third-party services like Stripe (payments) and AI bots to the core platform.
\end{itemize}

\par I did not stay in a single lane. If a feature required a database migration, I wrote it. If the build pipeline failed, I fixed it. This holistic approach allowed me to understand the software lifecycle in its entirety.

\section{The Reality Check: Academia vs. Industry}

\par There is a massive disconnect between what you learn in school and what you do at work. Here’s the truth:

\begin{itemize}
    \item \textbf{Theory vs. Reality:} School teaches you the syntax; work teaches you the context. Knowing *how* to write a loop is easy. Knowing *where* to put it so it doesn't crash the browser is harder.
    \item \textbf{Tooling Shock:} In university, you maybe use Git once. In the industry, if you don't know Git, Jira, or how to read a PR, you're useless.
    \item \textbf{The Human Factor:} School projects are usually solo or with friends. Real work is collaborating with people who might disagree with your code style.
    \item \textbf{Messy Constraints:} In class, you have all semester. In work, you have until Friday, and the specifications just changed this morning.
    \item \textbf{Legacy Code:} School start projects from scratch (file -> new). Work involves staring at 5-year-old code and trying to figure out what the original developer was thinking.
\end{itemize}

\section{Pain Points \& Growth}

\subsection{The Hard Stuff (Challenges)}
\begin{itemize}
    \item Jumping into a massive codebase is terrifying. You feel like you know nothing.
    \item Balancing "make it perfect" with "get it done" is a constant struggle.
    \item Refactoring old components without breaking everything else is stressful.
    \item Keeping up with the sprint pace—it’s fast.
    \item CSS media queries are finicky.
\end{itemize}

\subsection{Where I Grew}
\begin{itemize}
    \item My React/Next.js skills are night-and-day better.
    \item I learned how to actually manage my time.
    \item I'm comfortable with the tools now—VS Code and I are best friends.
    \item I learned how to take feedback on my code without taking it personally.
    \item I can actually debug things now instead of just staring at the screen.
\end{itemize}

\section{Detailed Plan}

\par I was at \textbf{Tanbel Inc.} from \textbf{October 13, 2025} to \textbf{December 5, 2025}. I didn't just sit around; I was busy.

\subsection{Resources}

\begin{enumerate}
    \item \textbf{People:} MD Tanvir Ahmed (my supervisor) and the other devs were lifesavers. Code reviews were where the real learning happened.
    \item \textbf{Technical Resources:} Access to development tools including Visual Studio Code, Git/GitHub, Chrome DevTools, Jira for task management, and the company's internal component libraries and design systems.
\end{enumerate}

\subsection{Timeline}

\par Here is exactly what I did, week by week. It was a busy two months.

\begin{table}[H]
\centering
\caption{Project Timeline Planning - Weekly Task Distribution}
\label{tab:training_timeline}
\resizebox{\textwidth}{!}{%
\begin{tabular}{|p{5.5cm}|c|c|c|c|c|c|c|c|}
\hline
\textbf{Tasks} & \textbf{Wk 1} & \textbf{Wk 2} & \textbf{Wk 3} & \textbf{Wk 4} & \textbf{Wk 5} & \textbf{Wk 6} & \textbf{Wk 7} & \textbf{Wk 8} \\
\hline
Env Setup \& CI/CD Optimization & \checkmark & & & & & & & \\
\hline
Mobile App Navigation \& UI Fixes & \checkmark & & & & & & & \\
\hline
Chatbot Service Deployment & & \checkmark & & & & & & \\
\hline
Staging Database \& Env Config & & \checkmark & & & & & & \\
\hline
Invoice \& Booking Logic Fixes & & & \checkmark & & & & & \\
\hline
CSV Export & Hauler Details Impl. & & & \checkmark & & & & \\
\hline
ArrayQ Site Deployment & Branding & & & & \checkmark & & & \\
\hline
Blog Page \& Dynamic Content Caching & & & & & \checkmark & & & \\
\hline
Blog Fixes \& App Store Review & & & & & & \checkmark & & \\
\hline
Mobile App Pricing & Discount Fixes & & & & & & \checkmark & \\
\hline
Stripe Connect \& Subscription Features & & & & & & & & \checkmark \\
\hline
Final Testing \& Documentation & & & & & & & & \checkmark \\
\hline
\end{tabular}%
}
\end{table}

\section{Risk Analysis: What Could Go Wrong?}

\begin{itemize}
    \item \textbf{Technical Challenges:} Sometimes I got stuck. The fix? Ask for help, read the docs, try again.
    \item \textbf{Timeline Risks:} Sprints move fast. I had to break tasks down so I didn't drown.
    \item \textbf{Communication Barriers:} Remote work is hard. Slack and Google Meet were the lifelines.
    \item \textbf{Quality:} We tested everything. I didn't want to be the guy who broke production.
\end{itemize}

\section{Report Outline}

\begin{itemize}
    \item \textbf{Chapter 2:} Who is \textbf{WeCycle}? The organization structure and mission.
    \item \textbf{Chapter 3:} What I actually did—tools and workflows.
    \item \textbf{Chapter 4:} The skills—what I learned.
    \item \textbf{Chapter 5:} Final thoughts—the conclusion.
    \item \textbf{Chapter 6:} Mapping outcomes (The official stuff).
    \item \textbf{Chapter 7:} SDG Goals (Sustainability mapping).
    \item \textbf{Chapter 8:} Appendix—The weekly logs.
\end{itemize}

\section{Conclusion}

\par This wasn't just an internship. It was a reality check. WeCycle didn't hold my hand; they let me work. I wrote real code for real users, and that changes how you think. It’s one thing to pass an exam. It’s another to fix a bug while your project manager is waiting. Honestly? I learned more here in a few weeks than I did in years of lectures. I’m not just ready to work. I’m ready to solve problems.